\documentclass[a4paper,10pt]{article}

\usepackage[utf8]{inputenc}
\usepackage{tabularx}
\usepackage[french]{babel}
\usepackage{csquotes}
\usepackage[none]{hyphenat}
%\usepackage[french=nohyphenation]{hyphsubst}
\usepackage{graphicx}
\usepackage{hyperref}
\usepackage{fancyhdr}
\usepackage{float}
\usepackage{caption}
\usepackage{listings}
\usepackage{spverbatim}
\usepackage{xcolor}
\usepackage[backend=bibtex]{biblatex}
\usepackage[]{algorithm2e}
\usepackage[mathscr]{euscript}
\usepackage{amsmath}

\captionsetup{justification=centering}

\lstdefinestyle{customc}{
  belowcaptionskip=1\baselineskip,
  breaklines=true,
  frame=L,
  xleftmargin=\parindent,
  language=C,
  showstringspaces=false,
  basicstyle=\footnotesize\ttfamily,
  keywordstyle=\bfseries\color{green!40!black},
  commentstyle=\itshape\color{purple!40!black},
  identifierstyle=\color{blue},
  stringstyle=\color{orange},
}

\lstdefinestyle{customasm}{
  belowcaptionskip=1\baselineskip,
  frame=L,
  xleftmargin=\parindent,
  language=[x86masm]Assembler,
  basicstyle=\footnotesize\ttfamily,
  commentstyle=\itshape\color{purple!40!black},
}

\lstdefinestyle{custombash}{
  belowcaptionskip=1\baselineskip,
  breaklines=true,
  frame=L,
  xleftmargin=\parindent,
  language=bash,
  showstringspaces=false,
  basicstyle=\footnotesize\ttfamily,
}

\lstset{escapechar=@}

\lhead{BAUDVIN CASSAGNES LAVILLE PUISSANT}
\rhead{Projet cryptographie}
%\rhead{2015-2016}

\pagestyle{fancy}

\sloppy

\addbibresource{biblio.bib}

\begin{document}

%\maketitle
\begin{titlepage}
      \begin{center}   
        \Huge
        \textbf{Projet de cryptographie}
        
        %\vspace{0.5cm}
        \LARGE
        ~
        
%         \vspace{0.5cm}
        
        \vfill
        \begin{figure}[H]
	    \centering
	    \begin{minipage}{0.9\textwidth}
		\centering
		\includegraphics[width=\textwidth]{./img/esiea.jpeg}
	    \end{minipage}\hfill
	\end{figure}
        \vfill
        
        \vspace{0.5cm}
        
        Implémentation et attaque d'un algorithme de chiffrement trappé
        
        \vspace{2cm}
        \textbf{Emmanuel Baudvin}\\\textbf{Lucien Cassagnes}\\\textbf{Lucas Laville}\\\textbf{Antoine Puissant}\\
        \vspace{0.8cm}
        \Large
        \underline{Enseignant :} M. Filiol\\
        \vspace{0.5cm}
        2015 - 2016%\today
        
    \end{center}
\end{titlepage}

\begin{abstract}
Le but de ce projet est de se familiariser avec les systèmes de chiffrements à trappes. Les deux systèmes de chiffrement \textit{SCEX\_T} qui nous sont présentés sont contiennent une trappe. Cela permet de les casser plus rapidement qu'en utilisant un méthode de type force brute. Ces deux cas sont inspirés de faits réels.\\
Le travail consiste à :
\begin{itemize}
 \item Implémenter les systèmes \textit{SCEX\_T1} et \textit{SCEX\_T2}.
 \item Analyser ces systèmes et identifier la trappe (analyse mathématique).
 \item Trouver une méthode afin d'exploiter cette trappe puis l'implémenter pour une attaque.
 \item Retrouver les textes clairs à partir des chiffrés fournis.
\end{itemize}
\end{abstract}

\newpage

\tableofcontents

\section{Implémentation des systèmes \textit{SCEX\_T1} et \textit{SCEX\_T2}}
\subsection{\textit{SCEX\_T1}}

\subsection{\textit{SCEX\_T2}}

\section{Analyse mathématique des systèmes \textit{SCEX\_T1} et \textit{SCEX\_T2}}
\subsection{\textit{SCEX\_T1}}
\subsubsection{Decimation cipher}
Ici, nous allons pouvoir retrouver une méthode de chiffrement dite de \enquote{\textit{décimation}}. Cette dernière se présente sous la forme suivante :
\begin{center}
    \begin{equation}
	c_t = \left\{
	    \begin{array}{lll}
		m_{t}\oplus\sigma_{t} & \mbox{si } t\equiv 0 & (\mbox{mod } 3) \\
		m_{t}\oplus\sigma_{t+1} & \mbox{si } t\equiv 1 & (\mbox{mod } 3) \\
		m_{t}\oplus\sigma_{t+2} & \mbox{si } t\equiv 2 & (\mbox{mod } 3)
	    \end{array}
	\right.
    \end{equation}
\end{center}

Le fonctionnement est très simple :

\begin{enumerate}
 \item Nous allons récupérer la position des caractères dans notre alphabet
 \item Nous allons alors multiplier/additionner/soustraire (modulo un nombre) cette position par un coefficient donné (la clé)
 \item La nouvelle position trouvée donne le chiffre du caractère clair
\end{enumerate}
Nous pouvons noter cet algorithme de la manière suivante :
\begin{center}
    \begin{equation}
	E(M)\equiv (kM+t)[n]
    \end{equation}
\end{center}
Où $k$ et $t$ sont des entiers et $k$ est premier par rapport à $n$ (le modulo utilisé, la taille de l'alphabet).\\
Pour déchiffrer un message ainsi codé, nous pouvons alors utiliser l'équation suivante :
\begin{center}
    \begin{equation}
	D(C)\equiv ((n-k)M+(n-t))[n]
    \end{equation}
\end{center}
Un des inconvénients de ce système est qu'il ne change pas la fréquence d'apparition des lettres, nous avons ici un système à permutations. Il est alors simple de reconnaître, en connaissant la langue d'origine du texte les fréquences les plus élevées. Une fois ces fréquences trouvées, nous pouvons poser deux équations à deux inconnues pour retrouver les coefficients $k$ et $t$.

\subsection{\textit{SCEX\_T2}}
\section{Exploitation et implémentation des trappes}
\subsection{\textit{SCEX\_T1}}

\subsection{\textit{SCEX\_T2}}
\section{Récupération des textes clairs}


\newpage
%\rhead{2015-2016}
\nocite{*}
\printbibliography\addcontentsline{toc}{section}{Références}

%\listoffigures\addcontentsline{toc}{section}{\listfigurename}
\end{document}
